\subsection{Reduced Norms and Determinants}

In this subsection, D is a field of finite degree over K with center K. We denote by \(D_{ab}^{*}\) the quotient of the multiplicative group \(D^{*}\) by its derived (or commutator) group and by \(\pi\) the canonical homomorphism from \(D^{*}\) to \(D_{ab}^{*}\) . The mapping \(Nrd_{D/K}\) induces a group homomorphism from \(D^{*}\) to \(K^{*}\) ; the kernel of this homomorphism contains the derived group of \(D^{*}\) because K is commutative. Hence there exists a unique homomorphism \(\overline{Nrd}\) from \(D_{ab}^{*}\) to \(K^{*}\) such that \(\mathrm{Nrd}_{\mathrm{D/K}}(x)=\overline{\mathrm{Nrd}}\circ\pi(x)\) for every \(x\in D^{*}\) .

PROPOSITION 9. --- Let V be a finite-dimensional right vector space over the field D. Let E be the algebra \(\operatorname{End}_{\mathrm{D}}(\mathrm{V})\) over the field K; it is central simple and of finite degree. For every invertible element u of E, we have

\[
\mathrm{Nrd} _ {\mathrm{E/K}} (u) = \overline {{\mathrm{Nrd}}} (\det u)\tag{64}
\]

(cf.~VIII, p.~452, Proposition 2 for the definition of the determinant \(\det u\) of \(u\) ).

We denote the dimension of V over D by n and identify E with the matrix algebra \(\mathbf{M}_{n}(\mathrm{D})\) using a basis of V over D. The multiplicative group \(\mathrm{GL}_{n}(\mathrm{D})\) of the algebra E is generated by the diagonal matrices and the matrices \(\mathrm{B}_{ij}(\lambda)\) (II, §10, No.~13, p.~362, Corollary 1). Proposition 9 therefore follows from the two specific cases below.

\begin{enumerate}
\def\labelenumi{\Alph{enumi})}
\tightlist
\item
  Suppose that u is the diagonal matrix \(\mathrm{diag}(a_{1},\ldots,a_{n})\) . For every \(1\leqslant i\leqslant n\) , let \(L_{i}\) be a maximal commutative subfield of D containing \(a_{i}\) ; let L be the subalgebra of E consisting of the diagonal matrices \(\mathrm{diag}(t_{1},\ldots,t_{n})\) with \(t_{i}\in L_{i}\) for \(1\leqslant i\leqslant n\) . Let d be the reduced degree of D over K. We have \([L_{i}:K]=d\) for \(1\leqslant i\leqslant n\) (VIII, p.~265, Corollary 2). The K-algebra L is isomorphic to \(L_{1}\times\cdots\times L_{n}\) and therefore semisimple of degree nd; now, we have \([E:K]=n^{2}[D:K]=n^{2}d^{2}=[L:K]^{2}\) . It follows that L is a maximal commutative semisimple subalgebra of E (VIII, p.~262, Proposition 3). By Proposition 7 of VIII, p.~346, we have \(\mathrm{Nrd}_{\mathrm{E}/\mathrm{K}}(u)=\mathrm{N}_{\mathrm{L}/\mathrm{K}}(u)\) . Consequently, by formula (18) of III, §9, No.~3, p.~544, we have
\end{enumerate}

\[
\begin{array}{c} \mathrm{Nrd} _ {\mathrm{E/K}} (u) = \prod_ {i = 1} ^ {n} \mathrm{N} _ {\mathrm{L} _ {i} / \mathrm{K}} (a _ {i}) = \prod_ {i = 1} ^ {n} \mathrm{Nrd} _ {\mathrm{D/K}} (a _ {i}) \\ = \mathrm{Nrd} _ {\mathrm{D/K}} (a _ {1} \dots a _ {n}) = \overline {{\mathrm{Nrd}}} (\pi (a _ {1} \dots a _ {n}))  . \end{array}
\]

Moreover, we have \(\det u = \pi(a_{1} \cdots a_{n})\) by Proposition 3 of VIII, p.~453, which gives formula (64) in this case.

\begin{enumerate}
\def\labelenumi{\Alph{enumi})}
\setcounter{enumi}{1}
\tightlist
\item
  Suppose that u is equal to \(\mathrm{B}_{ij}(\lambda)\) , where \(\lambda\) is an element of D and i,j are distinct integers in the interval \([1,n]\) . Denote by d the reduced degree of D over K and by M the vector space over K deduced from V by restriction of scalars from D to K. Then M is a simple E-module, and we have
\end{enumerate}

\[
\mathrm{Pc} _ {\mathrm{M/K}} (u; \mathrm{X}) = \mathrm{Pcrd} _ {\mathrm{E/K}} (u; \mathrm{X}) ^ {d}\tag{65}
\]

by formula (26) of VIII, p.~341. Moreover, M is a vector space of dimension \(nd^2\) over K, and \(u - 1_{\mathrm{M}}\) is a nilpotent endomorphism of M; we therefore have

\[
\mathrm{Pc} _ {\mathrm{M/K}} (u; \mathrm{X}) = (\mathrm{X} - 1) ^ {n d ^ {2}}.\tag{66}
\]

By comparing formulas (65) and (66), we obtain

\[
\mathrm{Pcrd} _ {\mathrm{E/K}} (u; \mathrm{X}) = (\mathrm{X} - 1) ^ {n d}\tag{67}
\]

and, in particular, \(\mathrm{Nrd}_{\mathrm{E}/\mathrm{K}}(u)=1\) . We also have det u=1 by Proposition 3 of VIII, p.~453; formula (64) therefore holds in this case.

Remark. --- We have \(\mathrm{Nrd}_{\mathrm{E}/\mathrm{K}}(u)=0\) if the element u of E is not invertible (VIII, p.~340, Proposition 3).

